\chapter{Formal Interface}
\label{ch:formal-interface}

\section{Real-line copies and regions}

\begin{definition}[Labeled real-line copies]
  \label{def:real-line-copies}
  \lean{Iut.RealLineCopy.Copy,Iut.RealLineCopy.Point,Iut.RealLineCopy.Transport}
  \leanok
  The formalization treats ordered real-line copies and positive-scale
  transports explicitly.  A comparison cannot silently identify a
  \ThetaPilot{} real line, a \qPilot{} real line, and the final signed
  log-volume target unless a transport has been supplied.
\end{definition}

\begin{definition}[Regions and measures]
  \label{def:regions-and-measures}
  \lean{Iut.RealLineCopy.Region,Iut.RealLineCopy.RegionComparison,Iut.RealLineCopy.RegionMeasure}
  \leanok
  Region comparisons package a source-to-target transport together with a
  target region.  Region measures are abstract monotone real-valued measures
  used to model log-volume estimates.
\end{definition}

\begin{definition}[Common target bound]
  \label{def:common-target-bound}
  \lean{Iut.RealLineCopy.RegionComparisonFamily.CommonTargetBound,Iut.RealLineCopy.TransportedRegionFamily}
  \leanok
  A common target bound records that every choice in a family of transported
  regions is bounded by the same measured target.  This is the formal shape of
  the post-hull upper-bound interface.
\end{definition}

\section{Algorithmic output and source ledger}

\begin{definition}[Algorithmic output]
  \label{def:algorithmic-output}
  \lean{Iut.RealLineCopy.AlgorithmicOutput}
  \leanok
  An algorithmic output packages transported region-family data together with
  qualitative certificate slots for \(\IPL\), \(\SHE\), \(\APT\), hull plus
  determinant, and common-container compatibility.
\end{definition}

\begin{definition}[Source obligation ledger]
  \label{def:source-obligation-ledger}
  \lean{Iut.Stage1.SourceObligationLedger}
  \leanok
  The source obligation ledger records the data required to promote a
  charted common-container comparison to the signed Corollary 3.12 endpoint.
  It keeps the chosen output, membership proof, target volume, \(\Theta\) bound,
  \qPilot{} positivity, and normalization proof in one auditable package.
\end{definition}

\begin{lemma}[Ledger three-term comparison]
  \label{lem:ledger-three-term-comparison}
  \uses{def:source-obligation-ledger}
  \lean{Iut.Stage1.SourceObligationLedger.qSigned_le_thetaSigned,Iut.Stage1.SourceObligationLedger.qSigned_le_thetaSigned_eq_threeTerm,Iut.Stage1.SourceObligationLedger.qSigned_le_thetaSigned_eq_membership_commonBound}
  \leanok
  The ledger proves \(q_{\mathrm{signed}}\le\Theta_{\mathrm{signed}}\) by the
  explicit chain
  \[
    q_{\mathrm{signed}}
      \le \text{chosen target volume}
      \le \Theta_{\mathrm{signed}}.
  \]
  The second comparison is tied to the charted common bound produced by the
  supplied certificate; it is not an arbitrary convenient \(\Theta\) bound.
\end{lemma}

\begin{proof}
  The Lean equalities expose the comparison as the ledger's three-term
  comparison and, more concretely, as membership in the chosen target followed
  by the common-bound inequality.
\end{proof}

\begin{definition}[Stage 1 comparison endpoint]
  \label{def:stage1-comparison-endpoint}
  \lean{Iut.Stage1.Stage1Comparison,Iut.Stage1.Corollary312ComparisonData}
  \leanok
  A Stage 1 comparison endpoint is the final signed real comparison
  \[
    q_{\mathrm{signed}} \le \Theta_{\mathrm{signed}},
  \]
  tagged by the \qPilot{} and \ThetaPilot{} sides.  The endpoint type is not a
  source proof; the source route must construct it from the ledger.
\end{definition}

\begin{lemma}[Projection warning]
  \label{lem:projection-warning}
  \lean{Iut.Stage1.corollary312_from_stage1_comparison}
  \leanok
  Once a completed endpoint has been supplied, the formal Corollary
  3.12-shaped inequality is a projection from that endpoint.
\end{lemma}

\begin{proof}
  The Lean theorem intentionally does no source mathematics.  It prevents us
  from confusing the endpoint package with a proof that the package is
  constructible from \IUT{} I--III.
\end{proof}

\section{Signed real algebra}

\begin{definition}[Corollary 3.12 predicate shape]
  \label{def:corollary-312-predicate-shape}
  \lean{Iut.Stage1.Corollary312Inequality}
  \leanok
  The bare predicate states the target signed pilot inequality for
  already-tagged \ThetaPilot{} and \qPilot{} log-volumes.  It is a proposition
  shape, not by itself a proof that source obligations construct the needed
  comparison.
\end{definition}

\begin{lemma}[Signed comparison gives the corollary shape]
  \label{lem:signed-comparison-gives-corollary}
  \uses{def:corollary-312-predicate-shape}
  \lean{Iut.Stage1.corollary312_of_signed_le}
  \leanok
  If \(q_{\mathrm{signed}}\le \Theta_{\mathrm{signed}}\), then the corresponding
  tagged pilot log-volumes satisfy the Corollary 3.12 inequality.
\end{lemma}

\begin{proof}
  This is a direct sign conversion: the stored pilot log-volume value is the
  positive magnitude, while the source display is written in signed form.
\end{proof}

\begin{theorem}[Final ordered-real calculation]
  \label{thm:final-real-calculation}
  \lean{Iut.Stage1.IUTStage1Corollary312SignedCThetaBound.cTheta_ge_neg_one}
  \leanok
  Suppose \(q_{\mathrm{signed}}\le \Theta_{\mathrm{signed}}\),
  \(0<-q_{\mathrm{signed}}\), and
  \[
    \Theta_{\mathrm{signed}}\le
      \CTheta\cdot(-q_{\mathrm{signed}}).
  \]
  Then \(-1\le \CTheta\).
\end{theorem}

\begin{proof}
  This is pure ordered-real algebra.  The hard source issue is therefore not
  this calculation, but the construction of the signed comparison consumed by
  it.
\end{proof}

\begin{theorem}[Strict and boundary alternatives]
  \label{thm:strict-boundary-alternatives}
  \lean{Iut.Stage1.IUTStage1Corollary312SignedCThetaBound.cTheta_eq_neg_one_or_gt_neg_one}
  \leanok
  Under the same hypotheses, either \(\CTheta=-1\) or \(-1<\CTheta\).  In the
  boundary case, the formal corridor also proves equality of the signed
  \qPilot{} and \ThetaPilot{} readings and negativity of the \ThetaPilot{}
  reading.
\end{theorem}

\begin{proof}
  The endpoint is a dichotomy for the final signed real inequality.  It should
  be read only after the source route has supplied the comparison
  \(q_{\mathrm{signed}}\le \Theta_{\mathrm{signed}}\).
\end{proof}
